\chapter{Fazit und Ausblick}
\label{chap:fazit}

Ziel war es, den bislang undokumentierten MDV-Prozess für Verbrauch und Reichweite bei Mercedes-AMG methodisch zu erfassen, mit BPMN 2.0 in Signavio zu modellieren und in die übergeordnete Prozessarchitektur einzuordnen. Hierzu wurde der Ablauf in Interviews mit Vertretern von drei Fachbereichen aufgenommen und als Modell mit 13 Prozessschritten abgebildet. Das Modell macht die Schritte, die Verantwortlichkeiten und die Entscheidungen zwischen den bekannten Eingangs- und Ausgangsgrößen sichtbar und schließt damit die in Kapitel~\ref{chap:problemkontext} beschriebene Lücke. Da es mit den direkten Prozessbeteiligten fachlich abgestimmt wurde, kann der Ablauf nun gegenüber Stakeholdern und Prüfinstanzen dargestellt werden. Bei der Einordnung wurde der MDV-Prozess keiner der Kategorien Management-, Kern- oder Unterstützungsprozess eigenständig zugeordnet, da er auf der vierten Ebene liegt. Seine Stellung ergibt sich stattdessen als Teilprozess des WLTP-Prozesses. Das bislang bei einzelnen Mitarbeitenden liegende Prozesswissen kann nach der Veröffentlichung in Signavio als gemeinsame Grundlage genutzt werden. Der MDV-Prozess lässt sich damit für Westeuropa mit BPMN 2.0 in Signavio erfassen, modellieren und in die Prozessarchitektur einordnen.
\newline
Die Ergebnisse unterliegen jedoch Grenzen. Das Modell bildet ausschließlich Westeuropa ab, da sich die Abläufe je Markt unterscheiden, sodass eine Übertragbarkeit auf weitere Märkte nicht belegt ist. Darüber hinaus beruht die Aufnahme auf Interviews mit Vertretern von drei Fachbereichen und gibt zunächst nur deren Perspektiven wieder. Diese Einschränkung wurde durch die fachliche Abstimmung des Modells verringert, jedoch nicht aufgehoben. Ob die Nachvollziehbarkeit gegenüber Prüfinstanzen gegeben ist, wurde nicht überprüft. Daneben ist die Einordnung vorläufig, da der WLTP-Gesamtprozess derzeit an Euro 7 angepasst wird (vgl. Abschnitt~\ref{sec:verortung}). Für die weitere Arbeit bietet es sich an, die Einordnung nach Abschluss der Euro-7-Anpassung und der Gremienfreigabe erneut zu prüfen und das Modell anschließend in Signavio zu veröffentlichen. Darüber hinaus sollte das Modell regelmäßig mit der Praxis abgeglichen werden, da sich Abläufe und regulatorische Anforderungen ändern.